\section{What would change the answer}\label{sec:future}

Several changes could alter the answer, and they divide into changes to the
question and changes within it.

\paragraph{A private information advantage.} Satellite coverage of the open ocean
or a destination model that resolves the three-quarters of the stock this pipeline
cannot would change the question rather than the answer, and would make the study a
test of the vendor's product. Within public data the most promising destination
feature is geometric: a Gulf cargo bound for Europe must exit north-east through the
Straits of Florida while Asia-bound cargoes route south, and the project's notes
record that a north-east crossing is strong evidence of a European arrival about
thirteen days ahead. The crossing is inside terrestrial range, so it is a coverage
question, and the live system's slot allocation was changed to keep departing
vessels subscribed through the Straits. That feature was not available on the
decade archives and is not tested here.

\paragraph{A fused capture-corrected supply series.} NOAA and GFW have opposite
coverage gradients (Section~\ref{sec:groundtruth}). A deduplicated union of the two,
matched on vessel, terminal and day, with the residual under-count carried as an
exposure offset calibrated to the EIA total, would recover the early years as
levels rather than only as changes. The acceptance test is that the fused series
reconciles with EIA within a few percent in every year.

\paragraph{A valid horizon test.} The withdrawn test of Section~\ref{sec:parta}
needs an unconditional arrival target: A1's conditional term composed with a
departure-process term for cargoes not yet loaded, scored against nulls of the same
unconditional quantity. The composition is a genuine modelling task and was parked
because its downstream consumer, the spread, showed no response to the input it
would sharpen.

\paragraph{The untested variants.} Within public data, the remaining untested
approaches are joint rather than marginal fits (spike-and-slab or partial least
squares over the collinear block), a fitted regime model in place of an observable
split, distributed lags of the signals, and the twenty-six signal keys the panel does
not use. Each has a weaker prior than the three mechanisms already tested, and each
faces a larger multiplicity problem against a residual that has never exceeded
\bMaxRsqCov{} of variance. Any of them should be accompanied by its own
pre-registration, its own multiplicity correction, and the untouched-holdout
standard used here.

\paragraph{A longer coverage-consistent sample.} The power floor falls as $1/n$.
Halving it requires roughly six further years on the weekly grid, or a daily grid
with a target that is not a near-random walk at that cadence. The live feed is the
only source that continues, and its coverage is a fraction of the archives', so the
seam of Section~\ref{sec:seam} will not close on its own.

\paragraph{Process.} Two cheap changes are recorded as open work: a coverage-regime
guard on the modelling panel so that a level break at a feed boundary fails a
check rather than a plot, and a third-party timestamp on every future
pre-registration entry.
